\section{Teoremas varios}
\subsection{Fórmula integral de Cauchy}

\begin{teorema} [Fórmula integral de Cauchy]
    Sea $\Omega \subset \mathbb{C}$ dominio, $f \in  \mathcal{H}(\Omega)$ y $D$ un disco tal que $\overline{D} \subset \Omega$. Entonces, para todo $z \in D$ se cumple que:
    \[
    f(z) = \frac{1}{2 \pi i} \int_{\partial D} \frac{f(w)}{w - z} d w.
    \]    
\end{teorema}

\begin{proof}
    Sea $z \in D$, como $D$ es abierto, entonces existe $\rho > 0$ tal que el disco $D(z, \rho) \subset D$. Luego $d (z, \partial D) \geq \rho \geq 0$. Llamemos
    \[
        g_z : \Omega \setminus \{ z \} \to \mathbb{C}
    \]
    \[
        w \mapsto \frac{f(w)}{w - z}
    \]
    Es claro que $g_z$ es holomorfa en $\Omega \setminus \{ z \}$. Veamos en primer lugar que 
    \[
        \int_{\partial D} g_z(w) dw = \int_{\partial D_r} g_z(w) dw
    \]
    siendo $D_r = D(z, r)$ un disco tal que $0 < r < d(z, \partial D)$.\\
    Sea $I$ la extensión de un radio de $D_r$ que une $\partial D_r$ con $\partial D$, y consideramos la curva
    \[
        \Gamma = I \cup \partial D \cup I^{-1} \cup (\partial D_r)^{-1}.
    \]
    Tenemos que $\Gamma \sim 0$ en $\Omega \setminus \{ z \}$, luego como $g_z \in \mathcal{H}(\Omega \setminus \{ z \})$ se cumple por el teorema de Cauchy que
    \[
        \int_{\Gamma} g_z(w) dw = 0.
    \]
    así
    \[
        0 = \int_{I} g_z(w) dw + \int_{\partial D} g_z(w) dw + \int_{I^{-1}} g_z(w) dw + \int_{(\partial D_r)^{-1}} g_z(w) dw
    \]
    \[
        = \int_{\partial D} g_z(w) dw - \int_{\partial D_r} g_z(w) dw
    \]
    Por tanto $\int_{\partial D} g_z(w) dw = \int_{\partial D_r} g_z(w) dw$ para todo $0 < r < d(z, \partial D)$.\\
    Como $f \in \mathcal{H}(\Omega)$, la función
    \[
        h(w) = \begin{cases}
            \frac{f(w) - f(z)}{w - z} & w \neq z \\
            f'(z) & w = z
        \end{cases}
    \]
    es continua en $\Omega$, entonces es acotada en $\overline{D}$.\\
    Sea $M = \max_{w \in \overline{D}} |h(w)|$. Entonces, 
    \[
        \int_{\partial D_r} \frac{1}{w - z} dw = \int_{0}^{2 \pi} \frac{1}{\gamma(t) - z} \gamma'(t) dt
    \]
    tomando como parametrización $\gamma (t) = z + r e^{i t}$, $t \in [0, 2 \pi]$. Entonces $\gamma'(t) = i r e^{i t}$ y
    \[
        \int_{\partial D_r} \frac{1}{w - z} dw = \int_{0}^{2 \pi} \frac{1}{r e^{i t}} i r e^{i t} dt = \int_{0}^{2 \pi} i dt = 2 \pi i.
    \]
    Por tanto,
    \[
        \left| \int_{\partial D} \frac{f(w)}{w - z} dw - 2 \pi i f(z) \right| = \left| \int_{\partial D_r} g_z(w) dw - 2 \pi i f(z) \right|
    \]
    \[
        = \left| \int_{\partial D_r} g_z(w) dw - f(z) \int_{\partial D_r} \frac{1}{w - z} dw \right| = \left| \int_{\partial D_r} \left( \frac{f(w)}{w - z} - \frac{f(z)}{w - z} \right) dw \right|
    \]
    \[
        = \left| \int_{\partial D_r} \frac{f(w) - f(z)}{w - z} dw \right| = \left| \int_{\partial D_r} h(w) dw \right| \leq \Long(\partial D_r) \max_{w \in \partial D_r} |h(w)| \leq 2 \pi r M \xrightarrow{r \to 0} 0.
    \]
    Luego 
    \[
        \int_{\partial D} \frac{f(w)}{w - z} dw = 2 \pi i f(z).
    \]
\end{proof}

El resultado anterior nos permite calcular integrales, entre otras muchísimas cosas.

\ejemplo{
    Vamos a calcular la integral
    \[
        \int_{C(i,1) } \frac{z^2}{z^2+1} dz.
    \]
    \[
        \int_{C(i,1) } \frac{z^2}{z^2+1} dz = \int_{C(i,1) } \frac{z^2}{(z-i)(z+i)} dz = \int_{C(i,1) } \left( \frac{z^2}{z+1i} \right) \frac{1}{z - i} dz.
    \]
    Llamemos $f(z) = \frac{z^2}{z+1i}$. Entonces, $f \in \mathcal{H}(\mathbb{C} \setminus \{-i\})$ y $\overline{D}(i,1) \subset \mathbb{C} \setminus \{-i\}$, siendo $\mathbb{C} \setminus \{-i\}$ un dominio. Por tanto, podemos aplicar la fórmula integral de Cauchy para obtener
    \[
        = \int_{C(i,1) } \frac{f(z)}{z - i} dz = 2 \pi i f(i) = 2 \pi i \frac{i^2}{i + 1i} = 2 \pi i \frac{-1}{2i} = -\pi.
    \]
}

\begin{teorema} [Teorema del valor medio de Gauss]
    Sea $\Omega \subset \mathbb{C}$ un dominio, $f \in \mathcal{H}(\Omega)$ y supongamos que $\overline{D}(a, r) \subset \Omega$. Entonces se cumple que
    \[
        f(a) = \frac{1}{2 \pi} \int_{0}^{2 \pi} f(a + r e^{i t}) dt
    \] 
\end{teorema}

\begin{proof}
    Por la fórmula integral de Cauchy, tenemos que
    \[
        f(a) = \frac{1}{2 \pi i} \int_{\partial D(a, r)} \frac{f(w)}{w - a} dw.
    \]
    Tomando como parametrización de $\partial D(a, r)$ la curva $\gamma(t) = a + r e^{i t}$, $t \in [0, 2 \pi]$, tenemos que $\gamma'(t) = i r e^{i t}$ y
    \[
        = \frac{1}{2 \pi i} \int_{0}^{2 \pi} \frac{f(a + r e^{i t})}{a + r e^{i t} - a} i r e^{i t} dt = \frac{1}{2 \pi i} \int_{0}^{2 \pi} \frac{f(a + r e^{i t})}{r e^{i t}} i r e^{i t} dt
    \]
    \[
        = \frac{1}{2 \pi i} \int_{0}^{2 \pi} f(a + r e^{i t}) i dt = \frac{1}{2 \pi} \int_{0}^{2 \pi} f(a + r e^{i t}) dt.
    \]
\end{proof}

\begin{corolario}
    Sea $\Omega \subset \mathbb{C}$ un dominio y $f \in \mathcal{H}(\Omega)$ una función no constante. 
    Entonces $|f|$ no alcanza su valor máximo en el interior de $\Omega$.
\end{corolario}

\begin{proof}
    Supongamos, por contradicción, que existe $z_0 \in \Omega$ tal que
    \[
        |f(z_0)| = \max_{z \in \Omega} |f(z)|.
    \]
    Veamos entonces que $f$ es constante en $\Omega$, lo que contradice la hipótesis inicial.\\
    Como $\Omega$ es abierto, existe $\delta > 0$ tal que $\overline{D}(z_0, \delta) \subset \Omega$. 
    Por el teorema del valor medio de Gauss, se tiene
    \[
        f(z_0) = \frac{1}{2\pi} \int_{0}^{2\pi} f(z_0 + \delta e^{it}) \, dt.
    \]
    Aplicando la desigualdad triangular,
    \[
        |f(z_0)| \leq \frac{1}{2\pi} \int_{0}^{2\pi} |f(z_0 + \delta e^{it})| \, dt.
    \]
    Observemos que la integral $\frac{1}{2\pi} \int_{0}^{2\pi} |f(z_0 + \delta e^{it})| \, dt$ es el promedio de los valores de $|f|$ en la circunferencia $\partial D(z_0, \delta) \subset \Omega$.\\
    Como $|f(z_0)|$ es el máximo de $|f|$ en $\Omega$, se cumple $|f(z_0 + \delta e^{it})| \le |f(z_0)|$ para todo $t \in [0,2\pi]$. Luego $|f(z_0)|$ es una cota superior del promedio, y por tanto
    \[
        |f(z_0)| \leq \frac{1}{2\pi} \int_{0}^{2\pi} |f(z_0 + \delta e^{it})| \, dt \leq |f(z_0)|.
    \] 
    Así la desigualdad anterior solo puede cumplirse si
    \[
        |f(z_0 + \delta e^{it})| = |f(z_0)| \quad \text{para todo } t,
    \]
    y además los números complejos $f(z_0 + \delta e^{it})$ tienen el mismo argumento que $f(z_0)$. 
    En consecuencia, $f$ es constante en la circunferencia $\partial D(z_0, \delta)$.

    Además, $\forall r \in (0, \delta)$, podemos aplicar el mismo razonamiento en el disco $D(z_0, r)$ para concluir que $f$ es constante en la circunferencia $\partial D(z_0, r)$ y, por tanto, en todo el disco $D(z_0, \delta)$. Luego por la conexidad de $\Omega$, se concluye que $f$ es constante en todo $\Omega$. 
    Esto contradice la hipótesis inicial, y por tanto $|f|$ no puede alcanzar su máximo en el interior de $\Omega$.
\end{proof}



El teorema de Taylor, nos asegura que si $f \in \mathcal{H}(\Omega)$, entonces $f$ es localmente la suma de una serie de potencias.

\begin{proposición}
    Sea $\gamma : [a,b] \to \mathbb{C}$ una curva $\mathcal{C}^1$ ó $\mathcal{C}^1$ a trozos y sea $(f_n)_{n \in \mathbb{N}}$ una sucesión de funciones continuas en $\gamma ^* = \gamma([a,b])$ tal que para cada $n \in \mathbb{N}$ 
    \[
        \exists M_n > 0 : |f_n(z)| \leq M_n \quad \forall z \in \gamma^*,
    \]
    y se cumple que $\sum_{n=0}^{\infty} f_n < \infty$. Entonces, 
    \[
        \int_{\gamma} \left( \sum_{n=0}^{\infty} f_n(z) \right) dz = \sum_{n=0}^{\infty} \left( \int_{\gamma} f_n(z) dz \right).
    \]
\end{proposición}

\begin{proof}
    Sea $\Omega \subset \mathbb{C}$ un dominio y $f \in \mathcal{H}(\Omega)$. Entonces para cada $z_0 \in \Omega$ existe $(c_n)_{n \in \mathbb{N}} \subset \mathbb{C}$ sucesión tal que
    \[
        f(z) = \sum_{n=0}^{\infty} c_n (z - z_0)^n \quad \forall z \in D(z_0, R) \subset \Omega.
    \]
    Además, para cada $n \in \mathbb{N} \cup \{0\}$ se cumple que
    \[
        c_n = \frac{1}{2 \pi i} \int_{\partial D(z_0, r)} \frac{f(z)}{(z-z_0)^{n+1}} dz \quad 0 < r < d(z_0, \partial \Omega).
    \]
\end{proof}

\subsection{Teorema de Taylor}

\begin{teorema}[Teorema de Taylor]
    Sea $\Omega \subset \mathbb{C}$ un abierto y $f \in \mathcal{H}(\Omega)$, entonces $\forall z_0 \in \Omega$ existe una unica sucesión $(c_n)_{n = 0}^{\infty} \subset \mathbb{C}$ tal que:
    \[
    f(z) = \sum_{n = 0}^{\infty} c_n (z - z_0)^n \quad 
    \forall z \in D(z_0, d(z_0, \partial \Omega))
    \]
    Además, para cada $n \in \mathbb{N} \cup \{0\}$:
    \[
    c_n =
    \frac{1}{2\pi i} 
    \int_{\partial D(z_0, r)} \frac{f(w)}{(w - z_0)^{n + 1}} \, dw
    \]
    Con radio $0 < r < d(z_0, \partial \Omega)$.
\end{teorema}

\begin{proof}
    Sea $z_0 \in \Omega$, como $\Omega$ es abierto $\exists \delta > 0 : D(z_0, \delta) \subset \Omega$, luego $d(z_0, \partial \Omega) \geq \delta > 0$, sea tambien $\rho = d(z_0, \partial \Omega)$.\\
    Tomemos entonces $r \in (0, \rho)$ y consideremos $\overline{D}_r = \overline{D}(z_0, r) \subset D(z_0, \rho) \subset \Omega$. Por la fórmula integral de Cauchy, tenemos que para todo $z \in D_r$:
    \[
        f(z) =
        \frac{1}{2 \pi i} \int_{\partial D_r} \frac{f(w)}{w - z} dw =^*
    \]
    Como $w \in \partial D_r$, entonces:
    \[
    \frac{1}{w-z} = 
    \frac{1}{w - z_0 + z_0 - z} = 
    \frac{1}{(w - z_0)(1 - \frac{z - z_0}{w - z_0})} =
    \frac{1}{w - z_0} \cdot \left( \sum_{n=0}^{\infty} \left( \frac{z - z_0}{w - z_0} \right)^n \right)
    \]
    Usando en el último paso que $\left| \frac{z - z_0}{w - z_0} \right| < \frac{r}{|w-z_0|} = \frac{r}{r} = 1$, luego la serie converge, por lo que:
    \[
    =^* \frac{1}{2 \pi i} \int_{\partial D_r} \frac{f(w)}{w-z_0} \left( \sum_{n=0}^{\infty} \left( \frac{z - z_0}{w - z_0} \right)^n \right) dw =
    \frac{1}{2 \pi i} \int_{\partial D_r} f(w) \left( \sum_{n=0}^{\infty} \frac{(z - z_0)^n}{(w - z_0)^{n+1}} \right) dw =
    \]
    \[
    = \frac{1}{2 \pi i} \int_{\partial D_r} \sum_{n=0}^{\infty} \left(\frac{f(w) (z - z_0)^n}{(w - z_0)^{n+1}}\right)  dw =^{**}
    \]
    Como para cada $n \in \mathbb{N} \cup \{0\}$, la función $w \mapsto \frac{f(w) (z - z_0)^n}{(w - z_0)^{n+1}}$ es continua en $\partial D_r$, y esta acotada por: \[
    \sup_{w \in \partial D_r} \left| \frac{f(w) (z - z_0)^n}{(w - z_0)^{n+1}} \right| \leq \frac{M}{r^{n+1}} |z - z_0|^n
    \] 
    definiendo $M \geq \max_{w \in \overline{D_r}} |f(w)| > 0$. Sabemos que $M$ existe y es cota porque $\overline{D_r}$ es compacto y $f$ es continua en $\Omega$. Así, la serie:
    \[
    \sum_{n=0}^{\infty} \frac{M}{r^{n+1}} |z - z_0|^n =
    \frac{M}{r} \sum_{n=0}^{\infty} \left| \frac{z - z_0}{r} \right|^n
    \underbrace{<}_{\left| \frac{z - z_0}{r} \right| < \frac{r}{r} = 1} \infty
    \]
    Luego por el ejercicio 4 de la hoja 6 podemos intercambiar el sumatorio y la integral, obteniendo:
    \[
    =^{**} \frac{1}{2 \pi i} \sum_{n=0}^{\infty} (z - z_0)^n \left(\int_{\partial D_r} \frac{f(w)}{(w - z_0)^{n+1}} dw \right) =
    \sum_{n=0}^{\infty} \left( \frac{1}{2 \pi i} \int_{\partial D_r} \frac{f(w)}{(w - z_0)^{n+1}} dw \right) (z - z_0)^n
    \]
    Así podemos llamar:
    \[
    c_n =
    \frac{1}{2 \pi i} \int_{\partial D_r} \frac{f(w)}{(w - z_0)^{n+1}} dw, \quad n \in \mathbb{N} \cup \{0\}
    \]
    deduciendo que:
    \[
    f(z) = \sum_{n=0}^{\infty} c_n (z - z_0)^n \quad \forall z \in D(z_0, r).
    \]
    Por tanto $f$ es infinitamente derivable en $D(z_0, r)$, y en particular en $z_0$, con derivadas $f^{(k)}(z_0) = k! c_k, \quad \forall k \in \mathbb{N} \cup \{0\}$, luego:
    \[
    c_n = \frac{f^{(n)}(z_0)}{n!} \quad \forall n \in \mathbb{N} \cup \{0\}
    \]
    Y finalmente como $c_n$ no depende de $r$, podemos tomar $r \to \rho = d(z_0, \partial \Omega)$, obteniendo la serie de potencias en $D(z_0, \rho)$:
    \[
    f(z) = \sum_{n=0}^{\infty} c_n (z - z_0)^n \quad \forall z \in D(z_0, \rho)
    \]
\end{proof}

Este resultado nos dice varias cosas:
\begin{itemize}
    \item Si $f \in \mathcal{H}(\Omega)$, entonces $f$ es localmente la suma de una serie de potencias.
    \item Si $f \in \mathcal{H}(\Omega)$, entonces $f$ es infinitamente derivable en $\Omega$.
    \item Si $f \in \mathcal{H}(\Omega)$, entonces $f' \in \mathcal{H}(\Omega)$.
    \item 
    \[
    \begin{cases}
        f \text{ es continua en } \Omega \text{ dominio}\\
        f \text{ tiene primitiva en } \Omega
    \end{cases}    
    \implies f \in \mathcal{H}(\Omega).
    \]
    \item Si $f \in \mathcal{H}(\Omega)$, entonces para cada $z_0 \in \Omega$ y $n \in \mathbb{N} \cup \{0\}$:
    \[
    f^{(n)}(z_0) = \frac{1}{2 \pi i} \int_{\partial D_r} \frac{f(w)}{(w - z_0)^{n+1}} dw \quad 0 < r < d(z_0, \partial \Omega).
    \]
\end{itemize}

Gracias a este resultado, podemos ampliar la formula integral de Cauchy a las derivadas de funciones holomorfas.

\begin{teorema}[Fórmula integral de Cauchy para las derivadas]
    Sea $\Omega \subset \mathbb{C}$ un dominio, $f \in \mathcal{H}(\Omega)$ y $D$ un disco tal que $\overline{D} \subset \Omega$. Entonces, para todo $z \in D$ y $n \in \mathbb{N} \cup \{0\}$ se cumple que:
    \[
    f^{(n)}(z) = \frac{n!}{2 \pi i} \int_{\partial D} \frac{f(w)}{(w - z)^{n+1}} dw.
    \]
\end{teorema}

\begin{proof}
    Sea $n \in \mathbb{N} \cup \{0\}$ y $z \in D$, tomamos $r > 0$ tal que $\overline{D(z, r)} \subset D$. Como la función $w \mapsto \frac{f(w)}{(w - z)^{n+1}}$ es holomorfa en $\Omega \setminus \{ z \}$ y además podemos ver que $\partial D$ es homótopa a $I \cup \partial D(z, r) \cup I^{-}$ en $\Omega \setminus \{ z \}$, siendo $I$ una curva que une $\partial D$ con $\partial D(z, r)$, entonces por el teorema de Cauchy para curvas homótopas tenemos que:
    \[
        \int_{\partial D} \frac{f(w)}{(w - z)^{n+1}} dw = 
        \int_{\partial D(z, r)} \frac{f(w)}{(w - z)^{n+1}} dw
    \]
    Luego:
    \[
    f^{(n)}(z) =
    \frac{n!}{2 \pi i} \int_{\partial D(z, r)} \frac{f(w)}{(w - z)^{n+1}} dw =
    \frac{n!}{2 \pi i} \int_{\partial D} \frac{f(w)}{(w - z)^{n+1}} dw
    \]
    \textbf{Nota: } Se puede usar el teorema de Cauchy para convexos en lugar del teorema de Cauchy para curvas homótopas.
\end{proof}

\textbf{Nota: } Si $\Omega \subset \mathbb{C}$ es un dominio y $z \in \Omega$,  $D(z,r) \subset \overline{D} \subset \Omega$ y $g \in \mathcal{H}(\Omega)$, entonces por el teorema anterior se cumple que:
\[
\int_{\partial D} g(w) dw =
\int_{\partial D(z,r)} g(w) dw
\] 

\subsection{Teorema de Liouville}
\begin{teorema} [Teorema de Liouville]
    Sea $f \in \mathcal{H}(\mathbb{C})$ y supongamos que $\exists M > 0$ tal que:
    \[
        |f(z)| \leq M \quad \forall z \in \mathbb{C}
    \]
    Entonces $f$ es constante en $\mathbb{C}$.
\end{teorema}


\begin{proof}
    Fijemos $z_0 \in \mathbb{C}$ cualquiera, entonces por el teorema de Taylor y teniendo en cuenta que $d(z_0, \partial \mathbb{C}) = +\infty$, tenemos que existe una sucesión $(c_n)_{n \in \mathbb{N}} \subset \mathbb{C}$ tal que:
    \[
        f(z) =
        \sum_{n=0}^{\infty} c_n (z - z_0)^n \quad \forall z \in \mathbb{C}
    \]
    \[
        c_n = \frac{1}{2 \pi i} \int_{\partial D(z_0, r)} \frac{f(w)}{(w - z_0)^{n+1}} \, dw \quad \forall r > 0, \forall n \in \mathbb{N} \cup \{0\}
    \]
    Sea $n \in \mathbb{N}$
    \[
        |c_n| \leq \left| \frac{1}{2 \pi i} \int_{\partial D(z_0, r)} \frac{f(w)}{(w - z_0)^{n+1}} \, dw \right| \leq \frac{1}{2 \pi} (2 \pi r) \cdot \sup_{w \in C(z_0, r)} \left| \frac{f(w)}{(w - z_0)^{n+1}} \right| 
    \]    
    \[
    \underbrace{\leq}_{|w-z_0|^{n+1}=r^{n+1}} r \cdot \frac{M}{r^{n+1}} = \frac{M}{r^n} \quad \forall r > 0 \xrightarrow{r \to \infty} 0
    \]
    puesto que $n \geq 1$. Luego $c_n = 0$ para todo $n \in \mathbb{N}$, y por tanto:
    \[
        f(z) = c_0 \quad \forall z \in \mathbb{C}
    \]
    o sea, $f$ es constante en $\mathbb{C}$.
\end{proof}

\subsection{Teorema fundamental del Álgebra}
\begin{teorema} [Teorema fundamental del Álgebra]
    Todo polinomio no constante con coeficientes complejos tiene al menos una raíz en $\mathbb{C}$.    
\end{teorema}

\begin{proof}
    Es aplicación directa del teorema de Lioulle. Sea $p(z) = a_n z^n + a_{n-1} z^{n-1} + \ldots + a_1 z + a_0$ un polinomio de grado $n \geq 1$ con coeficientes complejos. Si $p$ no se anulase en ningún punto de $\mathbb{C}$, entonces tendríamos que 
    \[
        \frac{1}{p(z)} \in \mathcal{H}(\mathbb{C})
    \]
    Veamos que $\frac{1}{p(z)}$ es acotada en $\mathbb{C}$.\\
    Como $n \geq 1$, se tiene que $\lim_{z \to \infty} p(z) = \infty$, luego $\lim_{z \to \infty} \frac{1}{p(z)} = 0$. Por tanto, existe $R > 0$ tal que:
    \[
        \left| \frac{1}{p(z)} \right| < 1 \quad \forall z \in \mathbb{C} \setminus \overline{D}(0, R)(1)
    \]
    Como $\frac{1}{p(z)}$ es holomorfa en $\mathbb{C}$, es continua en $\mathbb{C}$, luego es acotada en el compacto $\overline{D}(0, R)$(2), luego por (1) y (2) se deduce que $\frac{1}{p(z)}$ es acotada en $\mathbb{C}$.\\
    Ahora por el teorema de Lioulle, $\frac{1}{p(z)}$ es constante en $\mathbb{C}$, luego $p(z)$ es constante en $\mathbb{C}$, lo cual es una contradicción, porque habíamos supuesto que $p(z)$ era un polinomio de grado $n \geq 1$. 
\end{proof}

\subsection{Teorema de Morera}
\begin{teorema} [Teorema de Morera]
    Sea $\Omega \subset \mathbb{C}$ un dominio y sea $f : \Omega \to \mathbb{C}$ una función continua tal que:
    \[
        \int_{\partial \Delta} f(z) dz = 0 \quad \forall \Delta \subset \Omega
    \]
    Entonces, $f \in \mathcal{H}(\Omega)$.    
\end{teorema}

\begin{proof}
    Sea $z_0 \in \Omega$, como $\Omega$ es abierto, existe $R > 0$ tal que $D(z_0, R) \subset \Omega$.\\
    Para cada $z \in D(z_0, R)$, definimos:
    \[
        F(z) = \int_{[z_0, z]} f(w) dw
    \]
    Así, $F : D(z_0,R) \to \mathbb{C}$ y como $\int_{\partial \Delta} f(z) dz = 0$ para todo triángulo $\Delta \subset D(z_0, R)$, se deduce (de forma análoga a como hicimos en el teorema de existencia de primitivas) que $F$ es derivable en $D(z_0, R)$ y que $F'(z) = f(z)$ para todo $z \in D(z_0, R)$.\\
    Por el teorema de Taylor tenemos que $F \in \mathcal{H}(D(z_0, R)) \implies F' = f \in \mathcal{H}(D(z_0, R))$.\\
    Luego como $z_0 \in \Omega$ era arbitrario, se deduce que $f \in \mathcal{H}(\Omega)$.
\end{proof}

Dos resultados en los que usaremos el teorema de Morera son (para probar que una función es holomorfa) son el toerema de convergencia de Weierstrass y el teorema de prolongación analítica.

\begin{definición} [Convergencia uniforme sobre compactos]
    Sea $\Omega \subset \mathbb{C}$ abierto y sea $(f_n)_{n \in \mathbb{N}}$ una sucesión de funciones de $\Omega$ en $\mathbb{C}$.\\
    Decimos que $(f_n)_{n \in \mathbb{N}}$ converge uniformemente sobre compactos a una función $f : \Omega \to \mathbb{C}$ si para cada subconjunto compacto $K \subset \Omega$ se cumple que
    \[
        f_n \xrightarrow[n \to \infty]{} f \quad \text{uniforme en } K
    \]
    es decir, 
    \[
        \sup_{z \in K} |f_n(z) - f(z)| \xrightarrow[n \to \infty]{} 0.
    \]
\end{definición}

\subsection{Teorema de Weierstrass}
\begin{teorema} [Teorema de Weierstrass]
    Sea $\Omega \subset \mathbb{C}$ un dominio y sea $(f_n)_{n \in \mathbb{N}}$ una sucesión de funciones en $\mathcal{H}(\Omega)$ que converge uniformemente sobre compactos de $\Omega$ a una función $f : \Omega \to \mathbb{C}$. Entonces, $f \in \mathcal{H}(\Omega)$ y la sucesión $(f_n')_{n \in \mathbb{N}}$ converge uniformemente sobre compactos de $\Omega$ a $f'$.
\end{teorema}

\begin{proof}
    Sea $f_n \in \mathcal{H}(\Omega)$ para todo $n \in \mathbb{N}$. En particular, cada $f_n$ es continua en $\Omega$. Como
    \[
        \begin{cases}
            f_n \xrightarrow[n \to \infty]{} f, \text{ uniformemente sobre compactos de }\Omega\\
            f_n \text{ continua en } \Omega \quad \forall n \in \mathbb{N},
        \end{cases}
    \]
    se concluye que $f$ también es continua en $\Omega$.

    Veamos ahora que $\displaystyle \int_{\partial \Delta} f(z)\,dz = 0$ para todo triángulo $\Delta \subset \Omega$.

    Sea $\Delta \subset \Omega$ un triángulo cualquiera. Entonces,
    \[
        \left| \int_{\partial \Delta} f(z)\,dz - \int_{\partial \Delta} f_n(z)\,dz \right|
        = \left| \int_{\partial \Delta} (f(z) - f_n(z))\,dz \right|
        \leq \Long(\partial \Delta)\, \sup_{z \in \partial \Delta} |f(z) - f_n(z)|
        \xrightarrow{n \to \infty} 0,
    \]
    pues $\Delta$ es un compacto contenido en $\Omega$ y $f_n \to f$ uniformemente sobre compactos.

    Por lo tanto,
    \[
        \int_{\partial \Delta} f(z)\,dz = \lim_{n \to \infty} \int_{\partial \Delta} f_n(z)\,dz = 0,
    \]
    ya que, por el teorema de Cauchy para triángulos, $\int_{\partial \Delta} f_n(z)\,dz = 0$ al ser $f_n \in \mathcal{H}(\Omega)$.

    Aplicando el teorema de Morera, se deduce que $f \in \mathcal{H}(\Omega)$.

    \vspace{0.2cm}
    Para demostrar la segunda parte, utilizaremos la fórmula integral de Cauchy para las derivadas. Observemos que
    \[
        f_n'(z) \xrightarrow{n \to \infty} f'(z) \text{ uniformemente sobre compactos de } \Omega
        \iff
        f_n' \xrightarrow{n \to \infty} f' \text{ uniformemente sobre discos } \overline{D} \subset \Omega.
    \]

    Sea $D$ un disco tal que $\overline{D} \subset \Omega$. Queremos probar que
    \[
        \sup_{z \in \overline{D}} |f_n'(z) - f'(z)| \xrightarrow[n \to \infty]{} 0.
    \]

    Dado que $\overline{D} \subset \Omega$, por el Lema (1) del Tema 8 existe $\delta > 0$ tal que
    \[
        \overline{D} = \overline{D}(z_0, r) \subset D(z_0, r + \delta) \subset \overline{D}(z_0, r + \delta) \subset \Omega.
    \]
    Sea $D_1 := D(z_0, r + \tfrac{\delta}{2})$, de modo que $\overline{D} \subset D_1 \subset \overline{D_1} \subset \Omega$.

    Por la fórmula integral de Cauchy para las derivadas, se tiene:
    \[
        \sup_{z \in \overline{D}} |f_n'(z) - f'(z)|
        = \sup_{z \in \overline{D}} \left| \frac{1}{2\pi i} \int_{\partial D_1} \frac{f_n(w)}{(w - z)^2}\,dw - \frac{1}{2\pi i} \int_{\partial D_1} \frac{f(w)}{(w - z)^2}\,dw \right|
    \]
    \[
        \leq \sup_{z \in \overline{D}} \left| \int_{\partial D_1} \frac{f_n(w) - f(w)}{(w - z)^2}\,dw \right|
        \leq \sup_{z \in \overline{D}} \left( \Long(\partial D_1) \sup_{w \in \partial D_1} \left| \frac{f_n(w) - f(w)}{(w - z)^2} \right| \right)
    \]
    \[
        = \Long(\partial D_1) \sup_{w \in \partial D_1} \left( \sup_{z \in \overline{D}} \left|\frac{f_n(w) - f(w)}{(w - z)^2} \right| \right)
    \]

    Como $z \in \overline{D}$ y $w \in \partial D_1$, se cumple que $|w - z| \geq \tfrac{\delta}{2}$, por lo que
    \[
        |(w - z)^{-2}| \leq \frac{4}{\delta^2}.
    \]
    En consecuencia,
    \[
        \sup_{z \in \overline{D}} |f_n'(z) - f'(z)|
        \leq \Long(\partial D_1)\, \frac{4}{\delta^2} \sup_{w \in \overline{D_1}} |f_n(w) - f(w)|
        \xrightarrow[n \to \infty]{} 0,
    \]
    ya que $f_n \to f$ uniformemente sobre compactos de $\Omega$.

    Por lo tanto, $(f_n')$ converge uniformemente a $f'$ sobre compactos de $\Omega$, como queríamos demostrar.
\end{proof}

\subsection{Teorema de prolongación analítica}
\begin{teorema} [Teorema de prolongación analítica]
    Sea $\Omega \subset \mathbb{C}$ un dominio y $z_0 \in \Omega$, sea $f : \Omega \setminus \{ z_0 \} \to \mathbb{C}$ tal que $f \in \mathcal{H}(\Omega \setminus \{ z_0 \})$ y supongamos que 
    \[
        \lim_{z \to z_0} (z - z_0) f(z) = 0,
    \]
    entonces $f$ admite extensión holomorfa a $\Omega$.
\end{teorema}

\begin{proof}
    Supongamos en primer lugar que existe el límite
    \[
        \lim_{z \to z_0} f(z) = \alpha \in \mathbb{C}.
    \]
    De esta forma, definiendo $f(z_0) = \alpha$, tenemos que $f$ es continua en todo $\Omega$.\\
    Para probar que es holomorfa, usaremos el teorema de Morera. Comprobemos que 
    \[
        \int_{\partial \Delta} f(z) dz = 0 \quad \forall \Delta \subset \Omega. 
    \]
    Fijamos un triángulo $\Delta \subset \Omega$.\\
    Puede ocurrir uno de los tres casos:
    \begin{enumerate}
        \item $z_0 \notin \Delta$. 
        \item $z_0 \in \partial \Delta$.
        \item $z_0 \in \text{int}(\Delta)$.
    \end{enumerate}
    En el primer caso, $\Delta$ es un compacto contenido en $\Omega \setminus \{ z_0 \}$ (que es abierto), luego por el lema 1 del tema 8 tenemos que
    \[
        \exists \Omega_0 \subset \Omega \setminus \{ z_0 \} \text{ dominio tal que } \Delta \subset \Omega_0 \subset \overline{\Omega}_0 \subset \Omega \setminus \{ z_0 \}.
    \]
    \[
        f \in \mathcal{H}(\Omega \setminus \{ z_0 \}) \implies f \in \mathcal{H}(\Omega_0)
    \]
    y $\partial \Delta$ curva cerrada en $\Omega_0$, luego por el teorema de Cauchy para convexos deducimos
    \[
        \int_{\partial \Delta} f(z) dz = 0.
    \]
    En el segundo caso, puede ocurrir que 
    \begin{itemize}
        \item $z_0$ sea un vértice de $\Delta$.
        \item $z_0$ no sea un vértice de $\Delta$.
    \end{itemize}
    Si $\Delta$ es el triángulo con vértices $a, b, c$ y $z_0 = a$, dividimos $\Delta$ en tres triángulos:
    \[
        T_{\varepsilon}, T^1_{\varepsilon}, T^2_{\varepsilon}
    \]
    Como $z_1 \notin T^1_{\varepsilon}, z_0 \notin T^2_{\varepsilon}$, por el primer caso se cumple que
    \[
        \int_{\partial T^1_{\varepsilon}} f(z) dz = 0, \quad \int_{\partial T^2_{\varepsilon}} f(z) dz = 0.
    \]
    Además, considerando orientación positiva en cada triángulo, se tiene que
    \[
        \int_{\partial \Delta} f(z) dz =
        \int_{\partial T_{\varepsilon}} f(z) dz +
        \int_{\partial T^1_{\varepsilon}} f(z) dz +
        \int_{\partial T^2_{\varepsilon}} f(z) dz =
        \int_{\partial T_{\varepsilon}} f(z) dz.
    \]
    Luego,
    \[
        \left| \int_{\partial \Delta} f(z) dz \right| =
        \left| \int_{\partial T_{\varepsilon}} f(z) dz \right| \leq
        \Long(\partial T_{\varepsilon}) \sup_{z \in \partial T_{\varepsilon}} |f(z)| \xrightarrow[\varepsilon \to 0]{} 0,
    \]
    pues $\Long(\partial T_{\varepsilon}) \xrightarrow[\varepsilon \to 0]{} 0$ y $f$ está acotada en un entorno de $z_0$ (por ser continua en $z_0$).\\
    Así, se concluye que
    \[
        \int_{\partial \Delta} f(z) dz = 0.
    \]
    Si $z_0$ no es un vértice de $\Delta$, dividimos $\Delta$ en dos triángulos $\Delta_1, \Delta_2$ de forma que $z_0$ sea un vértice de ambos triángulos. Por el caso anterior, se cumple que
    \[
        \int_{\partial \Delta_1} f(z) dz = \int_{\partial \Delta_2} f(z) dz
    \]
    y como 
    \[
        \int_{\partial \Delta} f(z) dz =
        \int_{\partial \Delta_1} f(z) dz +
        \int_{\partial \Delta_2} f(z) dz
    \]
    deducimos que
    \[
        \int_{\partial \Delta} f(z) dz = 0.
    \]
    En el tercer caso, supongamos que $z_0 \in \text{int}(\Delta)$. Dividimos $\Delta$ en dos triángulos $\Delta_1, \Delta_2$ de forma que $z_0 \in \partial \Delta_1 \cap \partial \Delta_2$. Por el segundo caso, se cumple que
    \[
        \int_{\partial \Delta_1} f(z) dz = 0, \quad \int_{\partial \Delta_2} f(z) dz = 0.
    \]
    Además, considerando orientación positiva en cada triángulo, se tiene que
    \[
        \int_{\partial \Delta} f(z) dz =
        \int_{\partial \Delta_1} f(z) dz +
        \int_{\partial \Delta_2} f(z) dz = 0.
    \]
    Hemos demostrado el teorema bajo la suposición de que existe el límite $\lim_{z \to z_0} f(z) = \alpha \in \mathbb{C}$. Podemos definir la función $g : \Omega \to \mathbb{C}$ como
    \[
        g(z) =
        \begin{cases}
            f(z)(z - z_0), & z \in \Omega \setminus \{ z_0 \}\\
            0, & z = z_0
        \end{cases}
    \]
    y tenemos que $g \in \mathcal{H}(\Omega \setminus \{ z_0 \})$ y continua en $z_0$, luego como existe el límite $\lim_{z \to z_0} g(z) = 0$, se deduce que $g \in \mathcal{H}(\Omega)$. \\
    Por el teorema de Taylor,
    \[
        g(z) =
        \sum_{n=0}^{\infty} c_n (z - z_0)^n \quad \forall z \in D(z_0, d(z_0, \partial \Omega))
    \]
    Como $0=g(z_0) = c_0$, se tiene que
    \[
        g(z) = \sum_{n=1}^{\infty} c_n (z - z_0)^n = (z - z_0) \sum_{n=1}^{\infty} c_n (z - z_0)^{n-1} = (z - z_0) \sum_{n=0}^{\infty} c_{n+1} (z - z_0)^n
    \]
    Así, $z \in \Omega \setminus \{ z_0 \}$,
    \[
        g(z) = (z - z_0) f(z) = (z - z_0) \sum_{n=0}^{\infty} c_{n+1} (z - z_0)^n
    \]
    Luego, $f(z) = \sum_{n=0}^{\infty} c_{n+1} (z - z_0)^n$ para todo $z \in \Omega \setminus \{ z_0 \}$, y por tanto definiendo $f(z_0) = c_1$, tenemos
    \[
        f(z) = \sum_{n=0}^{\infty} c_{n+1} (z - z_0)^n \quad \forall z \in D(z_0, d(z_0, \partial \Omega))
    \]
    y por tanto $f \in \mathcal{H}(D(z_0, d(z_0, \partial \Omega)))$, con lo cual es derivable en $z_0$ y por tanto $f \in \mathcal{H}(\Omega)$.
\end{proof}

\begin{observación}
    \vspace{-2em}
    \begin{itemize}
        \item Si $\exists \lim_{z \to z_0} (z - z_0) f(z)$, entonces $\lim_{z \to z_0} (z - z_0) f(z) = 0$.
        \item Si $f$ es acotada en un entorno de $z_0$, entonces $\lim_{z \to z_0} (z - z_0) f(z) = 0$.
        \item Se puede probar fácilmente que si $\Omega \subset \mathbb{C}$ es un dominio y $f \in \mathcal{H}(\Omega \setminus \{z_1, \ldots, z_n \})$, y $\lim_{z \to z_k} (z - z_k) f(z) = 0$ para todo $k \in \{1, \ldots, n\}$, entonces $f$ admite extensión holomorfa a $\Omega$.
    \end{itemize}
\end{observación}

\begin{proposición}
    Sea $\Omega \subset \mathbb{C}$ simplemente conexo y sea $f : \Omega \to \mathbb{C}$ tal que $f \in \mathcal{H}(\Omega)$ y sin anularse en $\Omega$. Entonces, existe una determinación holomorfa de $\log(f)$ en $\Omega$, es decir, existe $g \in \mathcal{H}(\Omega)$ tal que:
    \[
        e^{g(z)} = f(z) \quad \forall z \in \Omega.
    \]
\end{proposición}

\begin{proof}
    Por el teorema de Taylor, sabemos que $f \in \mathcal{H}(\Omega) \implies f' \in \mathcal{H}(\Omega)$.\\
    Si $f \in \mathcal{H}(\Omega)$ y $f(z) \neq 0$ para todo $z \in \Omega$, entonces $\frac{1}{f} \in \mathcal{H}(\Omega)$.\\
    Asi tenemos que $\frac{f'}{f} \in \mathcal{H}(\Omega)$. Luego por el teorema de existencia de primitivas (al ser $\Omega$ simplemente conexo) existe $F \in \mathcal{H}(\Omega)$ tal que:
    \[
        F'(z) = \frac{f'(z)}{f(z)} \quad \forall z \in \Omega.
    \]
    Veamos que existe $c \in \mathbb{C}$ tal que la función $g(z) := F(z) + c$ cumple lo que queremos.\\
    Ciertamente tenemos que $g \in \mathcal{H}(\Omega)$. Considerando la función
    \[
        h(z) := e^{-F(z)} f(z) \quad \forall z \in \Omega,
    \]
    Tenemos que $h \in \mathcal{H}(\Omega)$ y $h(z) \neq 0$ para todo $z \in \Omega$. Además,
    \[
        h'(z) = -F'(z) e^{-F(z)} f(z) + e^{-F(z)} f'(z) = e^{-F(z)} \left( -\frac{f'(z)}{f(z)} f(z) + f'(z) \right) = 0 \quad \forall z \in \Omega.
    \]
    Luego $h$ es contante en $\Omega$. Sea $\beta \neq 0$ tal que $h(z) = \beta$ para todo $z \in \Omega$. 
    Si $\beta = |\beta|e^{i\theta_0}$ tenemos c un logaritmo de $\beta$, por ejemplo, $c = \log{|\beta|} + i \theta_0$. Así, $e^c = \beta$ y
    \[
        e^{g(z)} = e^{F(z) + c} = e^c e^{F(z)} = h(z) e^{F(z)} = f(z) \quad \forall z \in \Omega. 
    \]
\end{proof}

Como consecuencia en las hipótesis del resultado anterior tenemos que para cada $\alpha \in \mathbb{C}$, existe una determinación holomorfa de $f^{\alpha}$ en $\Omega$ (basta tomar $e^{\alpha g(z)}$), en particular de $\sqrt{f}, \sqrt[3]{f}$, etc.

\begin{observación}
    Un conjunto en $\mathbb{C}$ es simplemente conexo si su complementario en la esfera de Riemann es conexo.
\end{observación}


